% -----------------------------------------------------------------------------
% -----------------------------------------------------------------------------

\begin{exbox}[四点 Green 函数到截面]
前面的 LSZ 公式如果只停留在“给每条外腿乘一个逆传播子”，初学者很容易把它记成形式操作。最简单的检查是把 $\varphi^4$ 树级四点函数从头带一遍。相互作用仍取
\begin{equation*}
 \Lag_{\rm int}=-\frac{g_4}{4!\,\hbar c}\varphi^4,
 \qquad [g_4]=1,
\end{equation*}
而物理四动量满足 $p_i^2=m^2c^2$。在动量空间中，四点连通函数的最低阶奇异部分可以写成
\begin{equation}
 \begin{aligned}
 \widetilde G_{4,c}^{(1)}
 =&\,(2\pi\hbar)^4\delta^{(4)}(p_1+p_2-p_3-p_4)
 \left(-\frac{\ii g_4}{\hbar^2c^2}\right)\\
 &\times\prod_{a=1}^{4}
 \frac{\ii\hbar^3c}{p_a^2-m^2c^2+\ii0^+}.
 \end{aligned}
 \label{eq:phi4-fourpoint-poles-dp64}
\end{equation}
这个式子里其实混着三种不同信息：四维 $\delta$ 函数来自顶角位置积分；中间常数来自局域四价顶角；四个简单极点只是“从相互作用点传播到四个外部场插入点”的外腿传播。实验制备的是渐近单粒子，而不是带着这四个外腿传播子的场插入，因此 LSZ 要做的正是把最后一层剥掉。

对四条外腿逐一乘
\begin{equation*}
 \frac{p_a^2-m^2c^2}{\ii\hbar^3c}
\end{equation*}
并取 $p_a^2\to m^2c^2$，四个极点逐个相消，只留下
\begin{equation*}
 (2\pi\hbar)^4\delta^{(4)}(P_f-P_i)
 \left(-\frac{\ii g_4}{\hbar^2c^2}\right).
\end{equation*}
按上述约化振幅的统一定义，把纯粹来自 SI 场归一化的公共因子 $1/(\hbar^2c^2)$ 提出后，真正进入散射母式的是无量纲结果
\begin{equation}
 \boxed{\mathcal M_{\varphi\varphi\to\varphi\varphi}=-g_4.}
 \label{eq:phi4-reduced-amplitude-dp64}
\end{equation}
因此“图上四条外腿”与“物理振幅里的四个外粒子”不是同一件东西：前者在 Green 函数里带单粒子极点，后者经过 LSZ 后只保留在壳外态标签。

\eqnrefs{eq:phi4-reduced-amplitude-dp64} 代入二体截面母式后，等质量弹性散射有 $|\mathbf p_f|=|\mathbf p_i|$；末态两个标量完全相同，所以有额外的 $1/2!$。于是
\begin{equation*}
 \frac{\dd\sigma}{\dd\Omega}
 =\frac12\frac{\hbar^2}{64\pi^2s}g_4^2
 =\frac{\hbar^2g_4^2}{128\pi^2s},
\end{equation*}
这就是接触四点相互作用的各向同性二体截面。由于接触顶角没有任何 $t$ 或 $u$ 依赖，角分布各向同性；积分整个立体角后
\begin{equation}
 \boxed{
 \sigma_{\rm tot}
 =\frac{\hbar^2g_4^2}{32\pi s}.}
 \label{eq:phi4-total-cross-dp64}
\end{equation}
若用质心能量 $E_{\rm CM}=c\sqrt{s}$ 表示，便有
\begin{equation*}
 \sigma_{\rm tot}
 =\frac{g_4^2}{32\pi}
 \left(\frac{\hbar c}{E_{\rm CM}}\right)^2.
\end{equation*}
这最后一行也提供了一个快速量纲检查：四维无导数接触相互作用的无量纲耦合不提供长度尺度，因此高能截面只能按 $E_{\rm CM}^{-2}$ 衰减。作为数量级练习，若取 $g_4=0.50$、$E_{\rm CM}=1.00\,\mathrm{GeV}$，用 $\hbar c=0.19732698\,\mathrm{GeV\,fm}$ 得
\begin{align*}
 \sigma_{\rm tot}
 &=\frac{0.50^2}{32\pi}
 \left(\frac{0.19732698\,\mathrm{GeV\,fm}}{1.00\,\mathrm{GeV}}\right)^2,\\
 &\simeq9.68\times10^{-5}\,\mathrm{fm^2}
 =9.68\times10^{-35}\,\mathrm{m^2}
 \simeq0.968\,\mu\mathrm b.
\end{align*}
这个数值没有新的物理输入，只是在训练从自然出现的 $\hbar c/E$ 长度尺度回到实验常用的截面单位。


\medskip

\eqnrefs{eq:phi4-fourpoint-poles-dp64} 已把最低阶四点函数分成总四动量守恒、局域四价顶角和四条单粒子外腿极点。LSZ 约化作用在四条单粒子外腿极点上；直接检查四个在壳极点怎样逐个与约化算符相消。这里保留树级 $Z=1$，外部标量均视为稳定单粒子态。

对第 $a$ 条外腿定义逆极点因子
\[
 R_a(p_a)
 \equiv
 \frac{p_a^2-m^2c^2}{\ii\hbar^3c}.
\]
由\eqnrefs{eq:phi4-fourpoint-poles-dp64} 中的单腿因子可得
\[
 \lim_{p_a^2\to m^2c^2}
 R_a(p_a)
 \frac{\ii\hbar^3c}{p_a^2-m^2c^2+\ii0^+}=1.
\]
四条外腿分别取同一极限后，传播子极点全部消失，而顶角与总四动量守恒保留：
\[
 \lim_{p_a^2\to m^2c^2}
 \prod_{a=1}^{4}R_a(p_a)\,
 \widetilde G_{4,c}^{(1)}
 =(2\pi\hbar)^4\delta^{(4)}(P_f-P_i)
 \left(-\frac{\ii g_4}{\hbar^2c^2}\right).
\]
按前文约化振幅的 SI 归一化约定，把公共场归一化因子 $1/(\hbar^2c^2)$ 提出，得到
\[
 \ii\mathcal M=-\ii g_4,
 \qquad
 \mathcal M_{\varphi\varphi\to\varphi\varphi}=-g_4,
\]
即\eqnrefs{eq:phi4-reduced-amplitude-dp64}。


\medskip

\eqnrefs{eq:phi4-reduced-amplitude-dp64} 已经是可直接送入散射相空间的无量纲在壳振幅。剩下只需处理二体通量、全同末态计数和立体角积分；这一阶段不再涉及 LSZ 或 Green 函数。等质量弹性散射满足 $|\mathbf p_f|=|\mathbf p_i|$，而两个末态是同种实标量，因此二体截面母式中的对称因子为 $1/2!$。

由\eqnrefs{eq:2to2-cross-section-master-si-r9}，
\[
 \frac{\dd\sigma}{\dd\Omega}
 =\frac1{2!}\frac{\hbar^2}{64\pi^2s}
 |\mathcal M|^2
 =\frac{\hbar^2g_4^2}{128\pi^2s}.
\]
接触振幅没有 $t$ 或 $u$ 依赖，因此角分布各向同性。积分 $\int\dd\Omega=4\pi$ 得
\[
 \sigma_{\rm tot}
 =4\pi\frac{\hbar^2g_4^2}{128\pi^2s}
 =\frac{\hbar^2g_4^2}{32\pi s},
\]
即\eqnrefs{eq:phi4-total-cross-dp64}。
\end{exbox}




